\begin{abstract}
Language models (LMs) are becoming the foundation for almost all major language technologies, but their capabilities, limitations, and risks are not well understood.
We present Holistic Evaluation of Language Models (\benchmarkname) to improve the transparency of language models.
First, we taxonomize the vast space of potential scenarios (\ie use cases) and metrics (\ie desiderata) that are of interest for LMs.  
Then we select a broad subset based on coverage and feasibility, noting what's missing or underrepresented (\eg question answering for neglected English dialects, metrics for trustworthiness).
Second, we adopt a multi-metric approach:
We measure \nummetriccategories metrics (accuracy, calibration, robustness, fairness, bias, toxicity, and efficiency)
for \emph{each} of \numcorescenarios core scenarios
to the extent possible (87.5\% of the time), ensuring that metrics beyond accuracy don't fall to the wayside, and that trade-offs across models and metrics are clearly exposed.
We also perform \numcomponents targeted evaluations, based on \numcomponentscenarios targeted scenarios, to more deeply analyze specific aspects (\eg knowledge, reasoning, memorization/copyright, disinformation).
Third, we conduct a large-scale evaluation of \nummodels prominent language models (spanning open, limited-access, and closed models) on all \numscenarios scenarios, including \numnewscenarios scenarios that were not previously used in mainstream LM evaluation. 
Prior to \benchmarkname, 
models on average were evaluated on just 17.9\% of the \core \benchmarkname scenarios, with some prominent models not sharing a single scenario in common.
We improve this to 96.0\%: now all \nummodels models have been densely benchmarked on a set of core scenarios and metrics under standardized conditions.
Our evaluation surfaces \numfindings top-level findings concerning the interplay between different scenarios, metrics, and models.
For full transparency, we release all raw model prompts and completions publicly\footnote{\benchmarkUrl{}} for further analysis, as well as a general modular toolkit for easily adding new scenarios, models, metrics, and prompting strategies.\footnote{\githubUrl{}}
We intend for \benchmarkname to be a living benchmark for the community, continuously updated with new scenarios, metrics, and models.
\end{abstract}